\section{Relation to existing literature and exact scope}

The theorem-level conflict is with Campana--Demailly--Peternell
\cite[Theorem~2.2 and Corollary~2.3, p.~680]{CDP2020}.  Their hypotheses match
the candidate's asserted $H^1(X;\mathbb Z)=H^2(X;\mathbb Z)=0$, algebraic
dimension one, and holomorphic algebraic reduction.  Their conclusion gives
$c_3(X)\leq0$ and, for a threefold homeomorphic to $S^6$, algebraic dimension
zero.  The candidate instead asserts $c_3(X)=2$ and $a(X)=1$.  The conflict is
therefore exact if the candidate assertions are true; it is not resolved by
observing that the candidate proof does not invoke \cite{CDP2020}.

The first local issue isolated below occurs in the opening reduction of
\cite[\S7, p.~692]{CDP2020}.  The explicit nodal-cubic family in
\cref{prop:nodal-counterexample} refutes that reduction as a formal
sheaf-theoretic inference.  A second issue is the replacement of monodromy
coinvariants by the full fibre homology in
\cite[Lemma~3.2, pp.~84--85]{CDP1998}; CDP20 later imports the resulting
rank argument.  Either defect might in principle admit a different repair
under the global hypotheses.  No such repair is presently known here.

For the singular fibre, analytic duality is checked in two independent ways.
The proof in \cref{thm:cartier-duality} derives the result from smooth Serre
duality \cite[Theorem~9.37]{Lee2024}; Ramis--Ruget's dualizing-complex theorem
provides the singular analytic comparison
\cite[Theorems~1--2, pp.~81--84, and Proposition~2, p.~91]{RamisRuget1970}.
Relative duality for a proper analytic morphism is also cited in its precise
form as Ramis--Ruget--Verdier, Invent. Math. 13 (1971), 261--283
\cite{RamisRugetVerdier1971}.  Analytic cohomology and base change is
applied in the form of B\u{a}nic\u{a}--St\u{a}n\u{a}\c{s}il\u{a},
Chapter~III, \S3, Theorem~3.4 and Corollary~3.10, together with Grauert's
local-freeness/continuity theorem in Chapter~III, \S4, Theorem~4.12
\cite{BanicaStanasila1976}.  These locators cover all degrees used by the
structure-sheaf direct-image proof, not only its top-degree specialization.
The algebraic top-degree analogue in the Stacks Project, Tag~0EX2
\cite{StacksProject}, is a useful cross-check but is not substituted for the
analytic theorem.

The algebraic nearby-cycle formalism in SGA~7~II, Expos\'e XIII
\cite{SGA7II1973}, and the semistable/Raynaud framework in SGA~7~I,
Expos\'e IX \cite{SGA7I1972}, do not by themselves cover every step for the
present proper holomorphic map, which is not assumed algebraic.  The needed
analytic formalism is nevertheless available.  Dimca's Definition~4.2.1
constructs nearby cycles for complex analytic varieties; Proposition~4.2.2
identifies their stalks with Milnor-fibre cohomology; Proposition~4.2.11
commutes nearby cycles with proper analytic direct image; and
Theorem~2.3.26 gives proper base change, all for coefficient rings including
\(\mathbb Z\) \cite{Dimca2004}.  Kashiwara--Schapira
\cite[\S8.6]{KashiwaraSchapira1990} is an independent formal comparison.
Thus the general integral analytic framework is verified.  None of these
references by itself identifies the candidate's specialization image with
the full monodromy-invariant lattice.  That candidate-specific conclusion,
including saturation, is proved in
\cref{thm:p0-integral-specialization} for the reconstructed proper cusp
model, using the explicit local equations and a supported Gysin
calculation.  Equality with the separate cellular collapse map remains the
open part of obligation C3.

For compact-support Poincar\'e duality with local coefficients and its
projection formula we use Bredon \cite[Chapter~VI]{Bredon1997}, with
Hatcher \cite[\S3.H]{Hatcher2002} as a concise topological overview.  This
is the precise input behind the multiplicative identification
\(H^2(B,j_*\mathcal L)\simeq M_G\) used in the Leray calculation.

The algebraic reduction of a compact complex torus is used in the exact form
given by Birkenhake--Lange \cite[Chapter~2, \S6]{BirkenhakeLange1999}, with
the independent statement in Pareschi--Popa--Schnell
\cite[\S8, p.~2437]{PareschiPopaSchnell2017}.  The toric chart facts are
compared with Fulton \cite{Fulton1993}.  The exact classical-topology formula
\(\pi_1(X_\Sigma)\cong N/N_\Sigma\) is cited from Cox--Little--Schenck
\cite[Theorem~12.1.10]{CoxLittleSchenck2011}, with Danilov's original toric
treatment \cite[\S9]{Danilov1978} as an independent comparison.  All
assertions specific to the locally finite infinite fan and the candidate's
proper-discontinuity argument are proved directly.

Finally, the sphere-recognition endpoint uses Hurewicz and homological
Whitehead exactly as stated by Hatcher
\cite[Theorem~4.32 and Corollary~4.33]{Hatcher2002}, Smale's generalized
Poincar\'e theorem \cite[Theorem~A]{Smale1961}, and the computation of the
group of homotopy $6$-spheres by Kervaire--Milnor
\cite[opening table]{KervaireMilnor1963}.  Those results close the recognition
step after the candidate's fundamental group and integral homology are
established in \cref{thm:h2-vanishing,prop:geometric-pi1}; they are not used
to establish those inputs.
